\scanpage{23}
\mspage{19}
% archive sheet number: 22

supports \inserted{$\operatorname{supp}^{\circ}$} sont ``\uline{subordonnés}'' à un même drapeau $\delta\Phi\cup\delta\Psi\cup\partial I=T$. [Idem pour un ensemble quelconque de polygones, ou de parties constructibles de $\unsure{M|\mathcal M}$ (p.ex. des supports)]\deleted{\unsure{ceci}} -- quitte à subdiviser les intervalles polygonaux dans $\Phi,\Psi$ (ce qui ne change pas le $\operatorname{supp}^{\circ}$).

Toutes les opérations \inserted{usuelles} qu'on peut faire, sur une famille quelconque d'\deleted{$\operatorname{supp}^{\circ}$}éléments de $\Sigma_{\mathrm{cons}}$ mutuellement commensurables, se ramènent aux opérations $\cup,\cap,C$ dans un $\unsure{\mathscr P|\mathcal P}(\Phi_T)$ convenable.
(NB si $T$ de long. $n$, avec $n$ intervalles et $n+1$ sommets ($n\geq1$)
\[
\operatorname{card}\Phi_T=2n+1
\]
et
\[
\underbrace{\operatorname{Card}\mathscr P(\Phi_T)}_{\substack{\text{nb de supports}\\\text{subordonnés à }T}}=2^{2n+1}.
\]

\textcircled{2} Cas d'un $I$ où on ne se donne pas $\omega_0,\omega_1$. On doit le voir comme déduit du cas précédent $(I,\partial I)$, en \unsure{passant au}
